% ============================================================
\chapter{Applications : Finance, \'Economie, M\'et\'eorologie}
\label{ch:applications}
% ============================================================

\section{Exemple motivant}
Les outils d\'evelopp\'es dans ce cours s'appliquent directement \`a des probl\`emes concrets : gestion de portefeuille, pr\'evision macro\'economique, mod\'elisation climatique.

\section{Finance : Mod\'elisation de la volatilit\'e}\index{finance}

\subsection{Valeur \`a risque (VaR)}\index{VaR}

\begin{definition}[VaR]
La \textbf{Value at Risk} au niveau $\alpha$ est le quantile $\alpha$ de la distribution des pertes :
\[
P(L_t > \text{VaR}_\alpha) = \alpha.
\]
Avec un mod\`ele GARCH : $\text{VaR}_\alpha = -\hat\mu_t + \hat\sigma_t\,q_\alpha$, o\`u $q_\alpha$ est le quantile de la distribution des innovations.
\end{definition}

\subsection{Application}

\begin{codeblock}[title=\textbf{Python --- VaR avec GARCH}]
\begin{minted}{python}
import numpy as np
import pandas as pd
from arch import arch_model

# Simuler des rendements type financier
np.random.seed(42)
n = 2000
omega, alpha, beta = 0.01, 0.08, 0.90
sigma2 = np.zeros(n)
r = np.zeros(n)
sigma2[0] = omega / (1 - alpha - beta)
for t in range(1, n):
    sigma2[t] = omega + alpha * r[t-1]**2 + beta * sigma2[t-1]
    r[t] = np.sqrt(sigma2[t]) * np.random.standard_t(5)

# Estimer GARCH(1,1) avec innovations t de Student
am = arch_model(r * 100, vol='Garch', p=1, q=1, dist='t')
res = am.fit(disp='off')

# VaR 1%
forecasts = res.forecast(horizon=1)
sigma_next = np.sqrt(forecasts.variance.values[-1, 0])
from scipy.stats import t as t_dist
nu = res.params['nu']
VaR_01 = sigma_next * t_dist.ppf(0.01, nu)
print(f"VaR 1% = {VaR_01:.4f}%")
\end{minted}
\end{codeblock}

\section{\'Economie : Pr\'evision macro\'economique}\index{economie@\'economie}

\subsection{PIB et indicateurs avanc\'es}

\begin{methode}
\begin{enumerate}
\item Charger les s\'eries trimestrielles (PIB, consommation, investissement).
\item Tester la stationnarit\'e (ADF). Diff\'erencier si n\'ecessaire.
\item Estimer un VAR ou un ARIMA saisonnier.
\item \'Evaluer par validation hors \'echantillon.
\end{enumerate}
\end{methode}

\begin{codeblock}[title=\textbf{Python --- Pr\'evision macro}]
\begin{minted}{python}
import pandas as pd
import numpy as np
from statsmodels.tsa.api import VAR
from statsmodels.tsa.stattools import adfuller

# Exemple avec données simulées (remplacer par FRED/Eurostat)
np.random.seed(0)
n = 200
gdp = np.cumsum(np.random.normal(0.5, 1, n))      # trend + bruit
cpi = 0.3 * gdp + np.cumsum(np.random.normal(0, 0.5, n))

df = pd.DataFrame({'GDP_growth': np.diff(gdp),
                    'Inflation': np.diff(cpi)})

# VAR
model = VAR(df)
results = model.fit(ic='bic')
print(f"Ordre sélectionné: VAR({results.k_ar})")
print(results.summary())

# Prévision
forecast = results.forecast(df.values[-results.k_ar:], steps=8)
print("Prévision PIB (8 trimestres):", forecast[:, 0])
\end{minted}
\end{codeblock}

\section{M\'et\'eorologie : Temp\'eratures et saisonnalit\'e}\index{meteorologie@m\'et\'eorologie}

\begin{codeblock}[title=\textbf{Python --- S\'erie de temp\'eratures}]
\begin{minted}{python}
import numpy as np
import pandas as pd
import matplotlib.pyplot as plt
from statsmodels.tsa.seasonal import seasonal_decompose
from statsmodels.tsa.statespace.sarimax import SARIMAX

# Simuler des températures mensuelles
np.random.seed(42)
n_years = 30
t = np.arange(n_years * 12)
seasonal = 10 * np.sin(2 * np.pi * t / 12)
trend = 0.01 * t  # réchauffement
noise = np.random.normal(0, 2, len(t))
temp = 15 + trend + seasonal + noise

df = pd.DataFrame({'temp': temp},
    index=pd.date_range('1990-01', periods=len(t), freq='ME'))

# Décomposition
decomp = seasonal_decompose(df['temp'], model='additive', period=12)
decomp.plot()
plt.savefig('ch11_decomp.pdf')
plt.show()

# SARIMA
train = df.iloc[:-24]
test = df.iloc[-24:]
model = SARIMAX(train, order=(1,0,1),
                seasonal_order=(1,1,1,12))
results = model.fit(disp=False)
pred = results.forecast(steps=24)

rmse = np.sqrt(np.mean((test['temp'] - pred.values)**2))
print(f"RMSE prévision 24 mois: {rmse:.2f}°C")
\end{minted}
\end{codeblock}

\section{R\'esum\'e m\'ethodologique}

\begin{center}
\begin{tabular}{lll}
\toprule
\textbf{Domaine} & \textbf{Mod\`ele typique} & \textbf{Enjeu principal} \\
\midrule
Finance & GARCH, EGARCH & Volatilit\'e, queues lourdes \\
Macro\'economie & VAR, VECM, ARIMA & Causalit\'e, co\"int\'egration \\
M\'et\'eorologie & SARIMA, espace d'\'etats & Saisonnalit\'e, tendance \\
\'Epid\'emiologie & SIR + espace d'\'etats & Non-lin\'earit\'e \\
Signal & Analyse spectrale & P\'eriodicites cach\'ees \\
\bottomrule
\end{tabular}
\end{center}

\section{Exercices}

\begin{exercice}[$\star$]
Calculer la VaR \`a 5\% pour un GARCH(1,1) avec $\hat\sigma_{t+1}=2\%$ et innovations gaussiennes.
\end{exercice}

\begin{exercice}[$\star\star$]
Comparer la pr\'evision SARIMA et le lissage exponentiel de Holt--Winters sur les temp\'eratures mensuelles. Lequel est le plus pr\'ecis ?
\end{exercice}

\begin{exercice}[$\star\star\star$ -- Projet]
Choisir un jeu de donn\'ees r\'eel dans l'un des trois domaines. Appliquer la m\'ethodologie compl\`ete : exploration, mod\'elisation (au moins 3 mod\`eles), diagnostic, pr\'evision hors \'echantillon, comparaison. R\'ediger un rapport de 10 pages.
\end{exercice}

\begin{formules}
\begin{align*}
\text{VaR}_\alpha &= -\hat\mu_t + \hat\sigma_t\,q_\alpha \\
\text{RMSE} &= \sqrt{\frac{1}{H}\sum(X_{n+h}-\hat{X}_{n+h})^2} \\
\text{DM} &= \frac{\bar{d}}{\sqrt{\hat\sigma_d^2/H}}\dto\mathcal{N}(0,1)
\end{align*}
\end{formules}
